\section{SOLID}

    \par Os princípios SOLID, um acrônimo cunhado por Michael Feathers com base nos trabalhos de \cite{livro:martin:cleanarch}, representam cinco diretrizes cruciais para o design de software orientado a objetos. Estes princípios foram concebidos para auxiliar os desenvolvedores a construir sistemas que sejam mais fáceis de entender, manter e evoluir, aplicando-se a classes em qualquer camada da aplicação para se alcançar um melhor design, o que, por sua vez, contribui para uma base de código flexível e de fácil manutenção \cite{livro:joshi:2016}.

    \par A adesão a estes cinco preceitos - Princípio da Responsabilidade Única (Single Responsibility Principle - SRP), Princípio Aberto/Fechado (Open/Closed Principle - OCP), Princípio da Substituição de Liskov (Liskov Substitution Principle - LSP), Princípio da Segregação de Interface (Interface Segregation Principle - ISP) e Princípio da Inversão de Dependência (Dependency Inversion Principle - DIP) é fundamental para o desenvolvimento de software de alta qualidade. Visto que tais preceitos oferecem um vocabulário comum e um conjunto de heurísticas para a discussão e avaliação da arquitetura de software, direcionando as práticas de desenvolvimento para a minimização de dependências e o aprimoramento da coesão do código \cite{artigo:ramachandrappa:2024}.

    \par Conforme \cite{livro:martin:cleanarch}, o estudo dos princípios SOLID é essencial para a aplicação de uma arquitetura limpa, tendo isso em vista, pode-se traçar um paralelo de tais princípios com o objetivo do presente trabalho, onde o SRP e o DIP são os mais ligados na reestruturação, guiando a quebra dos componentes MVC em classes mais coesas, como Casos de Uso, e invertendo as dependências para que as regras de negócio não dependam mais do \textit{framework}. 
    
    \par Adicionalmente, ISP define os contratos limpos entre as novas camadas, enquanto o OCP e o LSP garantem que a arquitetura resultante permita a adição de novas funcionalidades e a troca de implementações sem alterar o núcleo do sistema, para uma abordagem detalhada de cada um desses princípios, os mesmos são descritos abaixo, com sua aplicação em exemplos práticos.